\documentclass{article}

\title{The Chain Rule}
\author{Ali}
\date{}

\begin{document}

\maketitle

\section{Introduction}

This report explains the chain rule in calculus. The chain rule is used to find the derivative of a function that has another function inside it.

\section{Main Points}

\begin{enumerate}
    \item A composite function is a function inside another function, like $f(g(x))$.
    \item To differentiate it, we first differentiate the outer function.
    \item Then we multiply by the derivative of the inner function.
\end{enumerate}

\section{Formula}

\[
\frac{d}{dx} f(g(x)) = f'(g(x)) \cdot g'(x)
\]

For example, if $y = (3x + 1)^2$, then
\[
\frac{dy}{dx} = 2(3x + 1) \cdot 3 = 6(3x + 1)
\]

\end{document}
